% !TEX root = ./main.tex

\chapter*{Remerciements}
\addcontentsline{toc}{chapter}{Remerciements}

Au terme de ce travail, je tiens à exprimer ma profonde gratitude.

Avant toute chose, je rends grâce à \textbf{Dieu}, Tout-Puissant et Miséricordieux, qui m'a accordé la force, la santé et la persévérance nécessaires à l'accomplissement de ce mémoire.

Je souhaite adresser mes plus sincères remerciements à \textbf{mon encadreur, Docteur LALAHARISON Hanjarivo}, pour sa disponibilité, ses précieux conseils et la confiance qu'il m'a accordée tout au long de ce travail. Sa rigueur et sa bienveillance ont été pour moi une source d'inspiration.

Ce mémoire n'aurait pas vu le jour sans le soutien infaillible de \textbf{ma femme}, qui m'a encouragé tout au long de ce parcours. Merci d'avoir cru en moi même quand je doutais moi-même. Je dédie ce travail à toi, ainsi qu'à \textbf{nos enfants}, dont les sourires et les câlins m'ont permis de tenir jusqu'au bout de ce projet.

J'adresse également ma gratitude à \textbf{mes beaux-parents}, ainsi qu'à \textbf{ma belle-grand-mère}, pour leur soutien moral et leurs prières, qui m'ont été d'un grand réconfort.

Je remercie chaleureusement \textbf{mes sœurs}, ainsi que \textbf{leurs maris}, pour leurs soutiens et leurs encouragements constants. J'avoue avoir eu besoin de leur sévérité et de leur affection pour surmonter les moments difficiles.

Un grand merci également à \textbf{mon beau-frère} ainsi qu'à \textbf{son épouse}, pour leur soutien et leur gentillesse qui m'ont profondément touché.

Je remercie chaleureusement mes amis, en particulier trois d'entre eux, pour leur amitié sincère. Vous avez toujours été comme des grands frères pour moi.

Enfin, je pense avec une profonde émotion à \textbf{mes défunts parents et ma défunte tante}, qui auraient été si fiers de ce jour. Ce mémoire est aussi le leur.

\vspace{1cm}
\begin{flushright}
\textit{Fait à Antananarivo, le 23 juillet 2026}
\end{flushright}